\documentclass[a4paper, 11pt]{article}

\usepackage[T1]{fontenc}
\usepackage{fix-cm}
\usepackage{amsmath}
\usepackage[scaled=0.85]{beramono}

\usepackage[polish]{babel}

\usepackage{etoolbox}
\preto{\enumerate}{\par\nobreak}

\usepackage[
  top=2.5cm, bottom=2.5cm,
  left=2.8cm, right=2.8cm,
  headheight=22pt
]{geometry}

\usepackage{tocloft}
\setlength{\cftbeforesecskip}{0.3em}
\setlength{\cftbeforesubsecskip}{0pt}

\usepackage{xcolor}
\definecolor{primary}{HTML}{1A1A1A}
\definecolor{accent}{HTML}{1A1A1A}
\definecolor{lightbg}{HTML}{F0F0F0}
\definecolor{codebg}{HTML}{EBEBEB}
\definecolor{codefg}{HTML}{1A1A1A}
\definecolor{rulegray}{HTML}{999999}

\usepackage[nobottomtitles*]{titlesec}
\titleformat{\section}
  {\color{primary}\large\bfseries}
  {\color{accent}\thesection.}{0.6em}{}[{\color{rulegray}\titlerule[0.6pt]}]
\titleformat{\subsection}
  {\color{primary}\normalsize\bfseries}
  {\color{accent}\thesubsection.}{0.5em}{}
\titlespacing{\section}{0pt}{1.4em}{0.7em}
\titlespacing{\subsection}{0pt}{1.0em}{0.4em}

\usepackage{fancyhdr}
\pagestyle{fancy}
\fancyhf{}
\fancyhead[L]{\small\color{primary}\textbf{Wizualizacja grafu planarnego}}
\fancyhead[R]{\small\color{rulegray}Dokumentacja implementacyjna -- Java}
\fancyfoot[C]{\small\color{rulegray}\thepage}

\usepackage{enumitem}
\setlist[itemize]{
  leftmargin=1.8em,
  itemsep=0.3em,
  topsep=0.4em
}

\usepackage{tcolorbox}
\tcbuselibrary{skins, breakable}

\newtcolorbox{cmdbox}{
  enhanced,
  colback=codebg,
  colframe=primary,
  arc=3pt,
  boxrule=0.8pt,
  left=8pt, right=8pt, top=5pt, bottom=5pt,
  fontupper=\ttfamily\small\color{codefg}
}

\newtcolorbox{filebox}{
  enhanced,
  colback=lightbg,
  colframe=rulegray,
  arc=3pt,
  boxrule=0.6pt,
  left=10pt, right=10pt, top=6pt, bottom=6pt,
  fontupper=\ttfamily\small\color{primary}
}

\usepackage[hidelinks]{hyperref}
\usepackage{needspace}
\usepackage{booktabs}
\usepackage{array}

\setlength{\parskip}{0.4em}
\setlength{\parindent}{0pt}
\setlength{\emergencystretch}{3em}

\begin{document}

\begin{titlepage}
  \centering
  \vspace*{4cm}

  {\fontsize{32}{38}\selectfont Wizualizacja grafu\\planarnego}

  \vspace{0.8cm}

  {\large\scshape\color{rulegray} Dokumentacja implementacyjna -- cz\k{e}\'s\'c Java}

  \vspace{1.2cm}

  {\small\scshape Wiktor Godula\\ Przemys\l{}aw Feliniak}

  \vfill

  {\small\color{rulegray} \today}

\end{titlepage}

\tableofcontents
\newpage

% ============================================================
\section{Wstęp}
% ============================================================

Dokument stanowi przewodnik po kodzie źródłowym aplikacji okienkowej
\texttt{Wizualizacja grafu planarnego} --- narzędzia napisanego w języku
Java (biblioteka Swing), pełniącego rolę graficznej nakładki na program
obliczeniowy wykonany w języku C. Dokumentacja opisuje podział systemu
na moduły, wewnętrzne struktury danych, architekturę interfejsu
użytkownika, przepływ sterowania, zasady komunikacji między komponentami
oraz zastosowane wzorce projektowe. Jej celem jest ułatwienie zrozumienia,
pielęgnacji i rozbudowy kodu.

Aplikacja Java nie implementuje własnych algorytmów obliczeniowych.
Cała logika matematyczna pozostaje po stronie zewnętrznego
programu \texttt{graf\_planarny}, uruchamianego jako osobny proces
przez klasę \texttt{CalculationEngine}. Aplikacja odpowiada wyłącznie
za prezentację grafu, interakcję z użytkownikiem oraz przekazywanie
danych do i z silnika C.

% ============================================================
\section{Struktura plików i podział na moduły}
% ============================================================

Kod źródłowy podzielony jest na dziesi\k{e}\'c modu\l{}\'ow, ka\.zdy reprezentowany
przez osobny plik \texttt{.java}. Jedynymi typami danych współdzielonymi
między modułami są \texttt{Graph}, \texttt{Vertex} i \texttt{Edge},
reprezentujące model grafu.

\begin{center}
\begin{tabular}{>{\ttfamily}l p{9.2cm}}
\toprule
\normalfont\textbf{Plik} & \textbf{Odpowiedzialno\'s\'c} \\
\midrule
Main.java            & Punkt wejścia aplikacji. Główne okno \texttt{JFrame}, koordynacja wszystkich komponentów GUI, obsługa zdarzeń wysokiego poziomu (wczytywanie, zapis, uruchamianie algorytmu). \\
\addlinespace
Graph.java           & Definicja struktury grafu: mapa wierzchołków \texttt{TreeMap<Integer, Vertex>}, lista krawędzi \texttt{List<Edge>}, obliczanie obwiedni (bounds). \\
\addlinespace
Vertex.java          & Pojedynczy wierzchołek: identyfikator \texttt{id} oraz współrzędne \texttt{x}, \texttt{y} z akcesorami. \\
\addlinespace
Edge.java            & Pojedyncza krawędź: nazwa, referencje do wierzchołków źródłowego i docelowego, waga. \\
\addlinespace
FileParser.java      & Operacje wejścia/wyjścia: wczytywanie grafu (tekstowe i binarne), zapis wyników (tekstowy i binarny), konwersja little-endian $\leftrightarrow$ big-endian. \\
\addlinespace
CalculationEngine.java & Uruchamianie zewnętrznego programu \texttt{graf\_planarny} (C) jako procesu potomnego, interpretacja kodów powrotu. \\
\addlinespace
GraphPanel.java      & Komponent \texttt{JPanel} odpowiedzialny za renderowanie grafu, obsługę myszy (zaznaczanie, przeciąganie wierzchołków, tryb przesuwania), skalowanie i autofit. \\
\addlinespace
AppToolPanel.java    & Panel narzędziowy (prawy): wybór algorytmu, liczba iteracji, format danych, zoom, edycja współrzędnych, widoczność etykiet i wag. \\
\addlinespace
AppMenuBar.java      & Pasek menu: opcje wczytywania grafu (tekstowo i binarnie), zapisu wyników (tekstowo/binarnie) oraz okno pomocy. \\
\addlinespace
AppStatusPanel.java  & Pasek statusu (dolny): wyświetlanie bieżącego stanu aplikacji, liczby wierzchołków i krawędzi oraz tymczasowych komunikatów o błędach. \\
\bottomrule
\end{tabular}
\end{center}

\subsection{Zależności między modułami}

Centralnym punktem zależności jest klasa \texttt{Main}, która posiada
referencje do wszystkich pozostałych komponentów i pełni rolę mediatora.
Pozostałe zależności są następujące:

\begin{itemize}
  \item \texttt{GraphPanel} otrzymuje referencje do \texttt{Graph} (model)
    oraz do \texttt{Main} (do komunikacji zwrotnej),
  \item \texttt{AppToolPanel} otrzymuje referencje do \texttt{Main}
    oraz \texttt{GraphPanel} (do bezpośredniego sterowania zoomem),
  \item \texttt{AppMenuBar} i \texttt{AppStatusPanel} komunikują się
    z \texttt{Main} przez wywołania zwrotne (listenery),
  \item \texttt{FileParser}, \texttt{CalculationEngine} oraz
    \texttt{Graph} są klasami pomocniczymi wywoływanymi przez \texttt{Main},
    nie posiadają własnych zależności do GUI,
  \item \texttt{Vertex} i \texttt{Edge} są prostymi strukturami danych
    wykorzystywanymi przez \texttt{Graph}, \texttt{GraphPanel}
    i \texttt{FileParser}.
\end{itemize}

W architekturze zastosowano wzorzec \textbf{Mediator} (rolę mediatora
pełni klasa \texttt{Main}) oraz \textbf{Observer} (słuchacze zdarzeń
Swing, np. \texttt{ActionListener} w menu i przyciskach). Dzięki takiej
strukturze komponenty GUI nie komunikują się ze sobą bezpośrednio,
lecz zawsze za pośrednictwem \texttt{Main}, co upraszcza zarządzanie
spójnością stanu interfejsu.

% ============================================================
\section{Struktury danych}
% ============================================================

Podstawowe struktury danych reprezentujące graf zdefiniowane są
w osobnych plikach \texttt{Vertex.java}, \texttt{Edge.java}
oraz \texttt{Graph.java}. W odróżnieniu od części C, gdzie dominują
tablice dynamiczne i ręczne zarządzanie pamięcią, w części Java
wykorzystywane są kolekcje z biblioteki standardowej
(\texttt{TreeMap}, \texttt{ArrayList}).

\subsection{Vertex --- wierzchołek grafu}

\begin{filebox}
public class Vertex \{\\
\quad int id;\\
\quad double x, y;\\
\\
\quad public Vertex(int id, double x, double y);\\
\quad public double getX();\\
\quad public double getY();\\
\quad public void setX(double x);\\
\quad public void setY(double y);\\
\}
\end{filebox}

Klasa przechowuje identyfikator wierzchołka (\texttt{id}) oraz jego
bieżące współrzędne $(x, y)$ w wirtualnym obszarze roboczym
$[0, 5000] \times [0, 5000]$. Pola są pakietowo dostępne (modyfikator
domyślny), co umożliwia szybki dostęp z \texttt{GraphPanel} (renderowanie)
bez narzutu wywołań getterów/setterów w pętli rysowania. Metody
\texttt{getX()} i \texttt{getY()} są zachowane dla klas zewnętrznych
względem pakietu.

\subsection{Edge --- krawędź grafu}

\begin{filebox}
public class Edge \{\\
\quad private String name;\\
\quad private Vertex source;\\
\quad private Vertex target;\\
\quad private double weight;\\
\\
\quad public Edge(String name, Vertex source,\\
\quad \qquad\qquad\; Vertex target, double weight);\\
\quad public String getName();\\
\quad public double getWeight();\\
\quad public Vertex getSource();\\
\quad public Vertex getTarget();\\
\}
\end{filebox}

Klasa przechowuje nazwę krawędzi, referencje do obiektów \texttt{Vertex}
(źródło i cel) oraz wagę (liczba zmiennoprzecinkowa podwójnej precyzji).
Referencje do wierzchołków umożliwiają bezpośredni dostęp do współrzędnych
podczas rysowania, bez konieczności wyszukiwania w mapie.

\subsection{Graph --- główna struktura grafu}

\begin{filebox}
public class Graph \{\\
\quad Map<Integer, Vertex> vertices = new TreeMap<>();\\
\quad List<Edge> edges = new ArrayList<>();\\
\quad private double minX, maxX, minY, maxY;\\
\\
\quad public void addVertex(int id, double x, double y);\\
\quad public void addEdge(String name, int sId, int tId,\\
\quad \qquad\qquad\qquad\quad\;\; double weight);\\
\quad public int getVertexCount();\\
\quad public int getEdgesCount();\\
\quad public Map<Integer, Vertex> getVertices();\\
\quad public void calculateBounds();\\
\quad public double getMinX();\\
\quad public double getMaxX();\\
\quad public double getMinY();\\
\quad public double getMaxY();\\
\quad public void clear();\\
\}
\end{filebox}

Wierzchołki przechowywane są w mapie \texttt{TreeMap<Integer, Vertex>},
co gwarantuje kolejność rosnącą według identyfikatora podczas iteracji
(np. przy zapisie wyników) i zapewnia średni czas dostępu $O(\log V)$.
Krawędzie przechowywane są w liście \texttt{ArrayList<Edge>}.

Metoda \texttt{addEdge} automatycznie wyszukuje wierzchołki źródłowy
i docelowy w mapie; jeśli oba istnieją, tworzy nowy obiekt \texttt{Edge}
i dodaje go do listy. W przeciwnym razie krawędź jest ignorowana
(co zapobiega niespójności danych).

Pola \texttt{minX}, \texttt{maxX}, \texttt{minY}, \texttt{maxY}
przechowują obwiednię grafu obliczaną przez \texttt{calculateBounds()}.
Są one wykorzystywane przez \texttt{GraphPanel} do automatycznego
skalowania (\texttt{autofit()}).

% ============================================================
\section{Moduł \texttt{FileParser} --- wejście/wyjście}
% ============================================================

Klasa \texttt{FileParser} zawiera całą logikę związaną z odczytem
i zapisem plików. Wszystkie publiczne metody operują na przekazanym
obiekcie \texttt{Graph}, modyfikując go przez jego metody publiczne.

\subsection{Metody prywatne --- wczytywanie}

Wczytywanie danych odbywa się przez trzy metody prywatne, wywoływane
w zależności od formatu pliku:

\begin{itemize}
  \item \texttt{readVertices(String path, Graph graph)} --- wczytuje
    plik tekstowy zawierający listę wierzchołków (pierwsza linia: liczba
    wierzchołków; kolejne: \texttt{id x y}). Używana do odczytu plików
    wynikowych generowanych przez silnik C lub przez funkcję zapisu.
    Format wejściowy używa kropki dziesiętnej (\texttt{Locale.US}).

  \item \texttt{readEdges(String path, Graph graph)} --- wczytuje
    plik tekstowy z listą krawędzi w formacie \texttt{nazwa zrodlo cel waga}.
    Jeśli napotkany wierzchołek nie istnieje jeszcze w grafie, jest
    automatycznie tworzony z domyślną pozycją $(0, 0)$. Metoda pomija
    linie niepasujące do formatu (zabezpieczenie \texttt{hasNextInt()}
    przed wyjątkiem).

  \item \texttt{readBinary(String path, Graph graph)} --- wczytuje
    binarny plik wynikowy w formacie little-endian (zgodnym z wyjściem
    programu C na platformie x86). Struktura pliku: 4 bajty --- liczba
    wierzchołków $V$ (\texttt{int}), następnie $V$ rekordów po 20 bajtów
    każdy (\texttt{int id}, \texttt{double x}, \texttt{double y}).
    Konwersja little-endian $\rightarrow$ big-endian realizowana jest
    przez wewnętrzne odwracanie kolejności bajtów w buforze.
    Po wczytaniu wywoływane jest \texttt{calculateBounds()}.
\end{itemize}

\subsection{Metody publiczne}

\subsubsection{loadFullGraph}

\begin{cmdbox}
public void loadFullGraph(String nodePath, String edgePath,\\
\quad\qquad\qquad\qquad\quad\; Graph graph, boolean binary)\\
\quad\qquad\qquad\qquad\quad\; throws Exception;
\end{cmdbox}

\textbf{Parametry:}
\begin{itemize}
  \item \texttt{nodePath} --- ścieżka do pliku z listą wierzchołków
    (wynikowego z silnika C); dla trybu binarnego traktowana również
    jako źródło danych binarnych,
  \item \texttt{edgePath} --- ścieżka do pliku z listą krawędzi,
  \item \texttt{graph} --- obiekt grafu do wypełnienia,
  \item \texttt{binary} --- jeśli \texttt{true}, wierzchołki wczytywane
    są z pliku binarnego (\texttt{readBinary}); jeśli \texttt{false},
    z tekstowego (\texttt{readVertices}).
\end{itemize}

\textbf{Działanie:} Wczytuje wierzchołki (z formatu tekstowego lub
binarnego), a następnie krawędzie (zawsze z pliku tekstowego
\texttt{edgePath}). Czyszczenie grafu (\texttt{graph.clear()}) odbywa się wewnątrz metod \texttt{readVertices()} i \texttt{readBinary()}, nie bezpośrednio w \texttt{loadFullGraph()}.

\textbf{Zwraca:} nic; w przypadku błędu rzuca \texttt{Exception}
z opisem problemu.

\subsubsection{saveToText}

\begin{cmdbox}
public void saveToText(String path, Graph graph) throws Exception;
\end{cmdbox}

Zapisuje współrzędne wszystkich wierzchołków do pliku tekstowego.
Pierwsza linia zawiera liczbę wierzchołków. Każda kolejna linia:
\texttt{id x y} z precyzją 4 miejsc po przecinku
(\texttt{printf("\%d \%.4f \%.4f\textbackslash n")}). Używa \texttt{Locale.US}
dla wymuszenia kropki dziesiętnej.

\subsubsection{saveToBinary}

\begin{cmdbox}
public void saveToBinary(String path, Graph graph) throws Exception;
\end{cmdbox}

Zapisuje współrzędne do pliku binarnego w formacie little-endian
(kompatybilnym z programem C). Używa \texttt{ByteBuffer} z kolejnością
bajtów \texttt{LITTLE\_ENDIAN}. Struktura: 4 bajty --- liczba
wierzchołków, następnie rekordy 20-bajtowe (\texttt{int id},
\texttt{double x}, \texttt{double y}).

% ============================================================
\section{Moduł \texttt{CalculationEngine} --- silnik obliczeniowy}
% ============================================================

Klasa \texttt{CalculationEngine} odpowiada za uruchomienie zewnętrznego
programu \texttt{graf\_planarny} (C) jako procesu potomnego oraz za
interpretację jego kodu powrotu.

\subsection{runGraphAlgorithm}

\begin{cmdbox}
public void runGraphAlgorithm(String inputPath,\\
\quad\qquad\qquad\qquad\qquad\; String algorithm,\\
\quad\qquad\qquad\qquad\qquad\; int iterations,\\
\quad\qquad\qquad\qquad\qquad\; boolean binary,\\
\quad\qquad\qquad\qquad\qquad\; Graph graph)\\
\quad\qquad\qquad\qquad\qquad\; throws Exception;
\end{cmdbox}

\textbf{Parametry:}
\begin{itemize}
  \item \texttt{inputPath} --- ścieżka do pliku wejściowego z listą krawędzi,
  \item \texttt{algorithm} --- nazwa algorytmu (\texttt{"fr"} lub
    \texttt{"tutte"}); przekazywana jako argument \texttt{-a},
  \item \texttt{iterations} --- liczba iteracji (dla FR); przekazywana
    jako argument \texttt{-n},
  \item \texttt{binary} --- jeśli \texttt{true}, dodaje flagę
    \texttt{-f bin} do polecenia,
  \item \texttt{graph} --- obiekt grafu, do którego zostaną wczytane
    wyniki po pomyślnym zakończeniu procesu.
\end{itemize}

\textbf{Działanie:}
\begin{enumerate}
  \item Ustala ścieżkę do pliku wykonywalnego: \texttt{core/graf\_planarny}
    (Linux/macOS) lub \texttt{core/graf\_planarny.exe} (Windows).
  \item Buduje listę argumentów: \texttt{-i <inputPath>},
    \texttt{-o wynik.txt}, \texttt{-a <algorithm>},
    \texttt{-n <iterations>} oraz opcjonalnie \texttt{-f bin}.
  \item Uruchamia proces przez \texttt{ProcessBuilder}.
  \item Oczekuje na zakończenie procesu (\texttt{waitFor()}).
  \item Jeśli kod powrotu to 0, wywołuje \texttt{FileParser.loadFullGraph}
    dla pliku \texttt{wynik.txt} i ścieżki wejściowej.
  \item Jeśli kod powrotu $\neq$ 0, rzuca \texttt{Exception}
    z komunikatem odpowiadającym kodowi błędu (mapowanie kodów
    realizowane jest wewnątrz klasy).
  \item Jeśli proces nie może zostać uruchomiony (\texttt{IOException}),
    rzuca \texttt{Exception} z komunikatem o braku pliku wykonywalnego.
\end{enumerate}

% ============================================================
\section{Moduł \texttt{GraphPanel} --- panel renderowania grafu}
% ============================================================

Klasa \texttt{GraphPanel} rozszerza \texttt{JPanel} i odpowiada za całą
wizualną reprezentację grafu oraz interakcję myszą.

\subsection{Obszar roboczy i układ współrzędnych}

Wirtualny obszar roboczy ma stały rozmiar $5000 \times 5000$ pikseli
(\texttt{WORK\_WIDTH}, \texttt{WORK\_HEIGHT}). Współrzędne wierzchołków
przechowywane w obiektach \texttt{Vertex} odnoszą się do tego wirtualnego
obszaru. Podczas renderowania stosowane jest przekształcenie afiniczne:

\begin{enumerate}
  \item \textbf{Translacja} --- środek wirtualnego obszaru roboczego
    jest wyrównywany ze środkiem widocznego panelu. Dodatkowo stosowane
    jest przesunięcie \texttt{offsetX}/\texttt{offsetY} (tryb pan).
  \item \textbf{Skalowanie} --- współczynnik \texttt{zoom} (domyślnie
    $0.5$, zakres od $\max(W_{\text{panelu}}/5000,\; H_{\text{panelu}}/5000)$
    do $2.0$) jest kontrolowany suwakiem w panelu narzędziowym.
\end{enumerate}

Transformacja realizowana jest przez wywołania \texttt{g2.translate(tx, ty)}
i \texttt{g2.scale(zoom, zoom)} przed pętlą rysowania. Dzięki temu
rysowanie wierzchołków i krawędzi odbywa się bezpośrednio na podstawie
współrzędnych wirtualnych.

\subsection{Obsługa myszy}

Wewnętrzna klasa anonimowa \texttt{MouseAdapter} obsługuje trzy rodzaje
interakcji:

\begin{itemize}
  \item \textbf{Kliknięcie (mousePressed)} --- jeśli tryb \texttt{panMode}
    jest aktywny, zapamiętuje pozycję myszy do późniejszego przesuwania.
    W przeciwnym razie przeszukuje wierzchołki w promieniu
    $2 \cdot V\_RAD$ (metoda \texttt{findVertex}) i ustawia znaleziony
    obiekt jako \texttt{selectedVertex}.

  \item \textbf{Przeciąganie (mouseDragged)} --- w trybie pan:
    aktualizuje \texttt{offsetX}/\texttt{offsetY} o wektor przesunięcia
    myszy i ogranicza przesunięcie do granic obszaru roboczego
    (\texttt{clampOffsets()}). W trybie edycji: przesuwa zaznaczony
    wierzchołek do pozycji kursora (z ograniczeniem do widocznego
    obszaru i do granic wirtualnego obszaru roboczego). Przekroczenie
    granic zgłaszane jest do \texttt{Main.notifyVertexOutOfBounds()}.
\end{itemize}

\subsection{Autofit}

Metoda \texttt{autofit()} oblicza obwiednię grafu, wyznacza skalę
tak, aby graf zajmował ok. 80\% widocznego obszaru, i odpowiednio
przeskalowuje współrzędne wszystkich wierzchołków. Po przeskalowaniu
graf jest centrowany w wirtualnym obszarze $5000 \times 5000$.

\subsection{Normalizacja do obszaru roboczego}

Metoda \texttt{normalizeToWorkspace()} skaluje i przesuwa współrzędne
wszystkich wierzchołków tak, aby zmieściły się w obszarze
$[12, 4988] \times [12, 4988]$ z zachowaniem proporcji. Wywoływana
jest warunkowo z metody \texttt{Main.updateUIState()} --- jeśli po wczytaniu pliku lub uruchomieniu algorytmu wierzchołki wykraczają poza granice obszaru roboczego, wyświetlany jest dialog potwierdzenia (\texttt{JOptionPane.YES\_NO\_OPTION}), a po zatwierdzeniu przez użytkownika wywoływana jest \texttt{normalizeToWorkspace()}.

\subsection{Renderowanie}

Metoda \texttt{paintComponent} wykonuje kolejno:
\begin{enumerate}
  \item Ograniczenie przesunięcia kamery (\texttt{clampOffsets()}).
  \item Obliczenie translacji. Zastosowanie \texttt{translate} i \texttt{scale}.
  \item Narysowanie tła obszaru roboczego (jasnoszary prostokąt
    $5000 \times 5000$) i jego ramki.
  \item Narysowanie wszystkich krawędzi jako linii (ciemnoszary).
    Opcjonalnie wyświetlenie wag w punkcie środkowym krawędzi
    (gdy \texttt{showWeights == true}).
  \item Narysowanie wszystkich wierzchołków jako okręgów o promieniu
    \texttt{V\_RAD = 12}. Wierzchołek zaznaczony (\texttt{selectedVertex})
    wypełniany jest kolorem białym, pozostałe --- czarnym.
    Opcjonalnie wyświetlenie etykiet (gdy \texttt{showLabels == true}).
\end{enumerate}

% ============================================================
\section{Pozostałe moduły GUI}
% ============================================================

\subsection{AppMenuBar}

Rozszerza \texttt{JMenuBar}. Tworzy dwa menu:

\begin{itemize}
  \item \textbf{Plik} --- zawiera opcje:
    \begin{itemize}
      \item \texttt{Wczytaj graf (Tekstowo)} --- wywołuje \texttt{parent.openFile(false)},
      \item \texttt{Wczytaj graf (Binarnie)} --- wywołuje \texttt{parent.openFile(true)},
      \item separator,
      \item \texttt{Zapisz wyniki (Tekstowo)} --- wywołuje
        \texttt{parent.saveAction(false)},
      \item \texttt{Zapisz wyniki (Binarnie)} --- wywołuje
        \texttt{parent.saveAction(true)}.
    \end{itemize}
  \item \textbf{Pomoc} --- zawiera opcję \texttt{Instrukcja i dokumentacja},
    wyświetlającą modalne okno dialogowe ze skróconą instrukcją obsługi
    i odesłaniem do pełnej dokumentacji.
\end{itemize}

\subsection{AppToolPanel}

Rozszerza \texttt{JPanel} z układem \texttt{BoxLayout} (oś Y).
Podzielony na cztery sekcje wizualne:

\begin{enumerate}
  \item \textbf{Ustawienia silnika} --- rozwijana lista algorytmów
    (\texttt{Fruchterman-Reingold} / \texttt{Tutte}), pole iteracji
    (nieaktywne dla Tutte), rozwijana lista formatu danych
    (\texttt{Tekstowy} / \texttt{Binarny}), przycisk \texttt{Uruchom}.
  \item \textbf{Widok} --- suwak zoom (zakres $5\%$--$200\%$ na sliderze,
    efektywne minimum ograniczane dynamicznie przez \texttt{GraphPanel.getMinZoom()},
    które zapobiega oddaleniu poniżej rozmiaru widocznego panelu),
    przycisk \texttt{Autofit}.
  \item \textbf{Edycja wierzchołka} --- pola tekstowe $X$ i $Y$,
    etykieta z ID zaznaczonego wierzchołka, przycisk \texttt{Zastosuj}.
  \item \textbf{Widoczność} --- pola wyboru: \texttt{Etykiety},
    \texttt{Wagi}, \texttt{Przesuwanie (Pan)}.
\end{enumerate}

Komunikacja z \texttt{Main} odbywa się przez referencję \texttt{parent},
a z \texttt{GraphPanel} przez referencję \texttt{graphPanel} (do
bezpośredniego ustawiania zoomu). Metody dostępowe
(\texttt{getSelectedAlgorithm()}, \texttt{isBinaryFormat()},
\texttt{getIterations()}, \texttt{setIterationsEnabled()},
\texttt{setVertexInfo()},
\texttt{getXText()}, \texttt{getYText()},
\texttt{setXText()}, \texttt{setYText()},
\texttt{updateZoomSlider()}) umożliwiają \texttt{Main}
pobieranie i ustawianie wartości kontrolek.

\subsection{AppStatusPanel}

Rozszerza \texttt{JPanel}. Wyświetla bieżący stan aplikacji w formacie:

\begin{filebox}
Stan: <komunikat> | Wierzchołki: <N> | Krawędzie: <M>
\end{filebox}

Metoda \texttt{notifyError} wyświetla tymczasowy komunikat o błędzie
(kolor czerwony), który po 3 sekundach (\texttt{Timer}) automatycznie
powraca do poprzedniego stanu. Mechanizm ten stosowany jest wyłącznie
dla ostrzeżeń o próbie przesunięcia wierzchołka poza obszar roboczy ---
wszystkie pozostałe błędy obsługiwane są przez modalne okna dialogowe
\texttt{JOptionPane} w klasie \texttt{Main}.

% ============================================================
\section{Moduł \texttt{Main} --- punkt wejścia i przepływ sterowania}
% ============================================================

Klasa \texttt{Main} rozszerza \texttt{JFrame} i pełni rolę:
\begin{itemize}
  \item punktu wejścia aplikacji (metoda \texttt{main}),
  \item głównego okna (tytuł: \texttt{"Wizualizacja grafu planarnego"}),
  \item mediatora między wszystkimi komponentami GUI,
  \item kontrolera przepływu sterowania.
\end{itemize}

\subsection{Metoda \texttt{main}}

\begin{cmdbox}
public static void main(String[] args);
\end{cmdbox}

Wywołuje \texttt{SwingUtilities.invokeLater()} z konstruktorem
\texttt{Main}, co zapewnia utworzenie GUI w wątku EDT (Event Dispatch
Thread), zgodnie z wymaganiami Swing.

\subsection{Konstruktor \texttt{Main()}}

Wykonuje inicjalizację wszystkich komponentów:
\begin{enumerate}
  \item Tworzy obiekty modelu (\texttt{Graph}), parsera (\texttt{FileParser})
    i silnika (\texttt{CalculationEngine}).
  \item Tworzy \texttt{GraphPanel} i przekazuje mu referencje do grafu
    oraz do siebie.
  \item Ustawia \texttt{AppMenuBar} jako pasek menu okna.
  \item Tworzy \texttt{AppToolPanel} (prawy panel narzędziowy)
    i \texttt{AppStatusPanel} (dolny pasek statusu).
  \item Rozmieszcza komponenty w układzie \texttt{BorderLayout}:
    \texttt{GraphPanel} w centrum, kontener \texttt{rightContainer}
    (zawierający przycisk zwijania \texttt{toggleBtn} i \texttt{AppToolPanel})
    po prawej stronie, pasek statusu na dole. Przycisk \texttt{toggleBtn}
    (symbole \texttt{\guillemotleft}/\texttt{\guillemotright}) przełącza widoczność panelu narzędziowego.
  \item Dezaktywuje panel narzędziowy do czasu wczytania pierwszego pliku.
  \item Wywołuje \texttt{graphPanel.autofit()} po wyrenderowaniu UI.
\end{enumerate}

\subsection{Główne metody sterujące}

\subsubsection{openFile}

\begin{cmdbox}
public void openFile(boolean binary);
\end{cmdbox}

Otwiera modalne okno wyboru pliku (\texttt{JFileChooser}). Po wyborze:
\begin{enumerate}
  \item Czyści graf (\texttt{graph.clear()}).
  \item W trybie binarnym wczytuje dane przez
    \texttt{parser.loadFullGraph(..., true)}.
  \item W trybie tekstowym:
    \begin{enumerate}
      \item Wywołuje prywatną metodę pomocniczą \texttt{prepareCleanedFile()},
        która tworzy plik tymczasowy pozbawiony komentarzy (linii zaczynających się
        od \texttt{\#}) i pustych linii. Plik tymczasowy oznaczany jest przez
        \texttt{deleteOnExit()}.
      \item Wywołuje przez refleksję prywatną metodę \texttt{readEdges}
        klasy \texttt{FileParser} dla oczyszczonego pliku.
    \end{enumerate}
  \item Aktywuje panel narzędziowy.
  \item Aktualizuje stan UI (\texttt{updateUIState()}).
\end{enumerate}
W przypadku błędu wyświetla okno dialogowe z komunikatem.

\subsubsection{runAlgorithm}

\begin{cmdbox}
public void runAlgorithm();
\end{cmdbox}

Wywoływana przez przycisk \texttt{Uruchom} w panelu narzędziowym.
Waliduje stan (czy wczytano plik, czy graf nie jest pusty), parsuje
parametry (algorytm, liczba iteracji, format), a następnie wywołuje
\texttt{engine.runGraphAlgorithm()}. Po pomyślnym zakończeniu odświeża
widok (\texttt{autofit()}) i pasek statusu. W przypadku błędu wyświetla
komunikat z prefiksem \texttt{"Błąd silnika: "}.

\subsubsection{applyCoords}

\begin{cmdbox}
public void applyCoords();
\end{cmdbox}

Wywoływana przez przycisk \texttt{Zastosuj} w panelu edycji wierzchołka.
Waliduje wprowadzone współrzędne: sprawdza, czy mieszczą się w obszarze
roboczym $[12, 4988]$, a następnie czy są skończonymi liczbami
rzeczywistymi (tj.\ nie są \texttt{NaN} ani nieskończone).
Następnie aktualizuje pozycję zaznaczonego wierzchołka i odświeża widok.
W przypadku błędu formatu lub przekroczenia zakresu wyświetla ostrzeżenie.

\subsubsection{saveAction}

\begin{cmdbox}
public void saveAction(boolean binary);
\end{cmdbox}

Otwiera modalne okno zapisu pliku (\texttt{JFileChooser}).
W zależności od flagi \texttt{binary} wywołuje \texttt{parser.saveToText()}
lub \texttt{parser.saveToBinary()}.

\subsubsection{updateUIState}

\begin{cmdbox}
private void updateUIState(String msg);
\end{cmdbox}

Wspólna metoda aktualizacji stanu całego interfejsu, wywoływana po każdej
operacji modyfikującej graf (wczytanie pliku, uruchomienie algorytmu).
Wykonuje kolejno:
\begin{enumerate}
  \item Aktualizuje pasek statusu (komunikat, liczba wierzchołków i krawędzi).
  \item Wywołuje \texttt{graphPanel.autofit()} w celu dopasowania widoku.
  \item Sprawdza, czy współrzędne wierzchołków mieszczą się w obszarze roboczym ($12$--$4988$). Jeśli nie, wyświetla modalne okno dialogowe (\texttt{JOptionPane.YES\_NO\_OPTION}) z pytaniem o automatyczne przeskalowanie.
  \item Po zatwierdzeniu wywołuje \texttt{graphPanel.normalizeToWorkspace()}.
  \item Synchronizuje suwak zoomu z bieżącą wartością.
\end{enumerate}

% ============================================================
\section{Zarządzanie pamięcią i zasobami}
% ============================================================

Zarządzanie pamięcią w części Java opiera się
na automatycznym odśmiecaniu (Garbage Collector). Niemniej program
przestrzega kilku zasad:

\begin{itemize}
  \item \textbf{Referencje} --- obiekty \texttt{Vertex} są współdzielone
    między mapą \texttt{Graph.vertices} a listą krawędzi
    (\texttt{Edge.source}, \texttt{Edge.target}). Dzięki temu zmiana
    współrzędnych wierzchołka jest natychmiast widoczna we wszystkich
    krawędziach, bez konieczności synchronizacji.
  \item \textbf{Czyszczenie grafu} --- metoda \texttt{Graph.clear()}
    usuwa wszystkie referencje z mapy i listy, co umożliwia GC
    zwolnienie pamięci po starych obiektach.
  \item \textbf{Pliki tymczasowe} --- pliki tworzone przez aplikację
    (np. przy czyszczeniu komentarzy) są oznaczane przez
    \texttt{deleteOnExit()}, co zapewnia ich usunięcie przy zamykaniu JVM
    nawet w przypadku nieoczekiwanego zakończenia programu.
  \item \textbf{Strumienie} --- \texttt{Scanner} i \texttt{PrintWriter}
    są używane w blokach \texttt{try-with-resources}, co gwarantuje
    automatyczne zamknięcie strumieni po zakończeniu operacji.
  \item \textbf{Procesy zewnętrzne} --- proces \texttt{graf\_planarny}
    jest uruchamiany przez \texttt{ProcessBuilder}; metoda
    \texttt{waitFor()} blokuje wątek do zakończenia procesu,
    zapobiegając wyciekom procesów.
\end{itemize}

% ============================================================
\section{Kompilacja i uruchamianie}
% ============================================================

\subsection{Wymagania}

\begin{itemize}
  \item Java SE 8 lub nowsza,
  \item skompilowany program \texttt{graf\_planarny} (C) umieszczony
    w podkatalogu \texttt{core/} katalogu roboczego aplikacji,
  \item system operacyjny z graficznym środowiskiem okienkowym,
  \item (opcjonalnie) narzędzie \texttt{make} do uproszczonej kompilacji.
\end{itemize}

\subsection{Metoda uproszczona (zalecana)}

W głównym katalogu projektu znajduje się plik \texttt{Makefile},
który automatyzuje proces kompilacji i uruchomienia. Aby zbudować
i wystartować aplikację, wystarczy wykonać:

\begin{cmdbox}
make
\end{cmdbox}

Polecenie to kompiluje wszystkie pliki \texttt{.java} i natychmiast
uruchamia GUI. Jeśli silnik C w folderze \texttt{core/} nie był jeszcze
budowany, należy uprzednio wykonać:

\begin{cmdbox}
cd core\\
make\\
cd ..
\end{cmdbox}

\subsection{Kompilacja ręczna}

W przypadku braku narzędzia \texttt{make} można skompilować projekt
bezpośrednio:

\begin{cmdbox}
\# Kompilacja wszystkich plików źródłowych:\\
javac -encoding UTF-8 gui/*.java
\end{cmdbox}

W wyniku powstają pliki \texttt{.class} w katalogu \texttt{gui/}.
Aplikację uruchamia się poleceniem:

\begin{cmdbox}
java -cp gui Main
\end{cmdbox}

\subsection{Struktura katalogów wdrożeniowych}

\begin{filebox}
katalog\_projektu/\\
\quad core/\\
\quad \quad graf\_planarny \quad\# (Linux/macOS)\\
\quad \quad graf\_planarny.exe \quad\# (Windows)\\
\quad gui/\\
\quad \quad *.java \quad\# pliki źródłowe\\
\quad \quad *.class \quad\# skompilowane klasy Java\\
\quad Makefile
\end{filebox}

Program C musi mieć prawo do wykonania (na systemach Unix:
\texttt{chmod +x core/graf\_planarny}). Aplikacja Java nie przeszukuje
innych ścieżek --- oczekuje pliku wykonywalnego dokładnie w podkatalogu
\texttt{core/} bieżącego katalogu roboczego.

\end{document}